\begin{abstract}
Class-Incremental Domain Adaptation (CIDA)~\citep{kundu2020cida} adapts a trained classifier to a new domain that also contains new classes.
It works without the source data and needs only one labelled image per new class.
Since the paper has no public code, we reimplement CIDA from scratch in PyTorch and first run it exactly as the paper describes it.
This version falls well short of the reported accuracy, and we trace the gap to four causes that each call for one change.
The changes re-estimate BatchNorm statistics on the target and start the domain projections at the exact identity.
They also train one pseudo-label loss on confident samples only and freeze the ImageNet backbone.
With these changes, \TODO{headline results, e.g.\ ALL and PRIV on Digits M$\to$U and Office-31, against the paper}.
We also show that the paper's $3\sigma$ rule for negative samples accepts almost every candidate in high dimensions.
\end{abstract}
